% 009 - Gloire à son nom
% Mode: 4/4 • Tonalité: A♭ majeur • Tempo: 95 (noire)
% Références bibliques: Jn 19:17-18; Ép 1:7; Col 1:20; Hé 9:22; 1 Pi 1:18-19
% -----------------------------------------------------------------------------
% Remarques:
% - Ce bloc est autonome et peut être inclus dans votre gabarit de hinário.
% - Les commandes \index[firstline]{...} supposent que vous avez défini un index
%   pour les premières lignes (ex.: via imakeidx). Adaptez si nécessaire.
% - Si vous utilise un environnement/hackage spécialisé (ex. songs, verse, etc.),
%   remplacez la structure simple ci-dessous par vos macros personnalisées.
% -----------------------------------------------------------------------------

\pdfminorversion=4
\documentclass[]{article}
\usepackage[utf8]{inputenc}
\usepackage{amssymb,latexsym,amsmath}
\usepackage[a5paper,top=2cm,bottom=2cm,left=2cm,right=2cm,marginparwidth=1.75cm]{geometry}
\usepackage{graphicx}
\usepackage[colorlinks=true, allcolors=blue]{hyperref}
\usepackage{helvet}
\renewcommand{\familydefault}{\sfdefault}

\begin{document}

\hypertarget{h009}{}
\section*{009. Gloire à son nom}
\addcontentsline{toc}{section}{009. Gloire à son nom}
\label{h:009}
\index{Gloire à son nom}

\vspace{-0.8em}
\noindent{\small\textit{Jn 19:17-18; Ép 1:7; Col 1:20; Hé 9:22; 1 Pi 1:18-19}}

\vspace{0.8em}
\medskip
\noindent\textbf{1.} Là sur la croix, Jésus fut cloué,\\
Pour mes péchés, je l'ai supplié.\\
Son sang versé m'a purifié.\\
Gloire à son nom saint !

\medskip
\indent\textbf{Gloire à son nom si saint !}\\
\indent\textbf{Gloire à son nom si saint !}\\
\indent\textbf{Son sang versé m'a purifié.}\\
\indent\textbf{Gloire à son nom !}

\medskip
\noindent\textbf{2.} Je fus sauvé, lavé de mon mal,\\
Jésus demeure en moi, c'est vital.\\
Là sur la croix, Il m'a pris, loyal.\\
Gloire à son nom saint !

\medskip
\noindent\textbf{3.} Quelle source précieuse et fidèle !\\
J'y suis entré, plein de joie réelle.\\
Jésus me sauve et me garde pur.\\
Gloire à son nom saint !

\end{document}
